% =====================================================================
% PROJETO DE ENGENHARIA DE SOFTWARE — AULA 3 (16/11)
% Desenhar as partes do sistema
% Caso 15 · Reserva de laboratórios da faculdade
% Os desenhos ficam em figuras/ e são incluídos com \input.
% =====================================================================

\documentclass[a4paper,12pt]{article}

\usepackage[utf8]{inputenc}
\usepackage[T1]{fontenc}
\usepackage[brazilian]{babel}
\usepackage{times}
\usepackage[a4paper,top=3cm,left=3cm,right=3cm,bottom=2.5cm]{geometry}
\usepackage{xltabular}
\usepackage{array}
\usepackage{amssymb}
\usepackage{titlesec}
\usepackage{enumitem}
\usepackage{tikz}
\usetikzlibrary{arrows.meta,positioning}
\usepackage{caption}

\setlength{\parindent}{0pt}
\setlength{\parskip}{6pt}
\renewcommand{\arraystretch}{1.3}
\setlist{itemsep=0pt,topsep=2pt}
\captionsetup{font=small,labelfont=bf,labelsep=period}

\titleformat{\section}{\fontsize{13}{15}\bfseries}{\thesection.}{1em}{}
\titleformat{\subsection}{\fontsize{12}{14}\bfseries}{\thesubsection.}{1em}{}

\newcolumntype{Y}{>{\raggedright\arraybackslash}X}

\begin{document}
\pagestyle{plain}

\begin{center}
  {\small\bfseries PROJETO DE ENGENHARIA DE SOFTWARE}\\[4pt]
  {\fontsize{16}{20}\selectfont\bfseries Aula 3 · 16/11 · Desenhar as partes do sistema}\\[6pt]
  {\large Caso 15 · Reserva de laboratórios da faculdade}
\end{center}

\vspace{6pt}

\begin{xltabular}{\linewidth}{|l|Y|}
  \hline
  \textbf{Professor} & Sérgio Santana \\ \hline
  \textbf{Disciplina} & Laboratório de Engenharia de Software \\ \hline
  \textbf{Integrante} & Wilia Francisco da Silva \\ \hline
  \textbf{Regra para testar} & Reservas de 1 a 4 horas, para até 40 alunos \\ \hline
\end{xltabular}

% ---------------------------------------------------------------------
\section{O sistema em 3 camadas}

O sistema segue o modelo cliente e servidor e é dividido em três
camadas. O professor usa o celular ou o computador (cliente); as regras
são verificadas no servidor; as informações ficam guardadas no banco de
dados. A Figura~\ref{fig:camadas} mostra as camadas e como elas
conversam.

\begin{center}
  % Arquitetura em 3 camadas — Caso 15 · Reserva de laboratórios
\begin{tikzpicture}[
  camada/.style={draw, thick, rounded corners=3pt, text width=12cm, inner sep=7pt, align=left},
  seta/.style={-{Stealth[length=3mm]}, thick},
  rot/.style={font=\small, fill=white, inner sep=2pt}
]
  \node[camada, fill=blue!6] (ap) {%
    \textbf{1. Camada de apresentação} \hfill \emph{Celular ou computador do usuário}\\[3pt]
    Tela de login · Agenda do laboratório · Nova reserva · Minhas reservas\\
    Mostra as informações e recebe o que o usuário digita ou toca.};

  \node[camada, fill=orange!8, below=1.9cm of ap] (rn) {%
    \textbf{2. Camada de regras de negócio} \hfill \emph{Servidor}\\[3pt]
    Login e perfis (professor e coordenação) · Regras RN01 a RN18 ·
    Bloqueio de choque de horário · Reservas de 1 a 4 horas e até 40 alunos
    (30 no Lab 202) · Cálculo do horário de término · Avisos por e-mail};

  \node[camada, fill=green!7, below=1.9cm of rn] (bd) {%
    \textbf{3. Camada de dados} \hfill \emph{Banco de dados}\\[3pt]
    Usuários · Laboratórios · Reservas};

  \draw[seta] ([xshift=-3cm]ap.south) -- node[rot]{pedido: ``reservar o Lab 101''} ([xshift=-3cm]rn.north);
  \draw[seta] ([xshift=3cm]rn.north) -- node[rot]{resposta: reserva aceita ou erro} ([xshift=3cm]ap.south);
  \draw[seta] ([xshift=-3cm]rn.south) -- node[rot]{consulta e gravação} ([xshift=-3cm]bd.north);
  \draw[seta] ([xshift=3cm]bd.north) -- node[rot]{dados guardados} ([xshift=3cm]rn.south);
\end{tikzpicture}

  \captionof{figure}{Arquitetura em 3 camadas do sistema de reserva de laboratórios.}
  \label{fig:camadas}
\end{center}

\begin{xltabular}{\linewidth}{|l|l|Y|}
  \hline
  \textbf{Camada} & \textbf{Onde fica} & \textbf{O que fica nela} \\ \hline
  Apresentação & Celular ou computador & As telas: login, agenda do laboratório, nova reserva e minhas reservas. Exemplo: o formulário de nova reserva. \\ \hline
  Regras de negócio & Servidor & As regras RN01 a RN18. Exemplo: verificar se a reserva tem de 1 a 4 horas, se cabe no laboratório e se não choca com outra; calcular o horário de término. \\ \hline
  Dados & Banco de dados & Usuários, laboratórios e reservas. Exemplo: a reserva confirmada fica gravada aqui. \\ \hline
\end{xltabular}

\textbf{Exemplo do caminho de uma reserva:}
\begin{enumerate}[leftmargin=1.2cm]
  \item O professor preenche a tela \emph{Nova reserva} no celular e toca em \emph{Confirmar reserva}.
  \item O pedido vai para o servidor, que verifica as regras: duração de 1 a 4 horas, até 40 alunos (30 no Lab 202), horário de funcionamento e choque com outras reservas.
  \item Se tudo estiver certo, o servidor grava a reserva no banco de dados e envia o aviso por e-mail.
  \item O servidor responde ao celular, que mostra \emph{Reserva confirmada}. Se alguma regra falhar, a tela mostra o motivo do erro.
\end{enumerate}

% ---------------------------------------------------------------------
\section{Esboço das 3 telas principais}

As três telas abaixo são da versão 2.0 e cobrem o que falta na versão
1.0: a agenda do laboratório, a consulta de disponibilidade antes de
reservar e a edição das reservas.

\noindent\begin{tabularx}{\linewidth}{|l|c|Y|}
  \hline
  \textbf{Tela} & \textbf{Histórias} & \textbf{O que a tela resolve} \\ \hline
  1 · Agenda do laboratório & H01 & Mostra os horários livres e ocupados de cada laboratório (RF06). \\ \hline
  2 · Nova reserva & H02, H01 & Reserva pelo celular e mostra se o laboratório está livre antes de confirmar (RF07). A duração é escolhida numa lista de 1 a 4 horas, o que evita erros de digitação. \\ \hline
  3 · Minhas reservas & H03, H04, H06 & Cancela e edita reservas (RF03, RF08) e mostra o aviso de confirmação (RF10). \\ \hline
\end{tabularx}

\begin{center}
  \begin{minipage}[t]{0.48\linewidth}
    \centering
    % Tela 1 · Agenda do laboratório (H01 · RF06)
\begin{tikzpicture}[font=\sffamily\scriptsize]
  % aparelho e barra superior
  \draw[thick, rounded corners=8pt] (0,0) rectangle (5.6,10.4);
  \fill[gray!25] (0.2,9.6) rectangle (5.4,10.2);
  \node[anchor=west] at (0.3,9.9) {\textbf{Reserva de Laboratórios}};
  \node[anchor=east] at (5.3,9.9) {Sair};
  \node[anchor=west] at (0.3,9.25) {\small\textbf{Agenda do laboratório}};

  % escolha do laboratório e do dia
  \draw (0.3,8.4) rectangle (5.3,8.9);
  \node[anchor=west] at (0.4,8.65) {Lab 101 (40 lugares)};
  \node[anchor=east] at (5.2,8.65) {$\blacktriangledown$};
  \node at (0.6,7.95) {$\blacktriangleleft$};
  \node at (2.8,7.95) {\textbf{Ter, 17/11/2026}};
  \node at (5.0,7.95) {$\blacktriangleright$};

  % horários do dia
  \draw (0.3,7.05) rectangle (5.3,7.5);
  \node[anchor=west] at (0.4,7.27) {07:00 \quad Livre — toque para reservar};
  \draw[fill=gray!30] (0.3,5.75) rectangle (5.3,7.05);
  \node[anchor=north west, text width=4.8cm] at (0.4,6.95) {\textbf{08:00 às 12:00} \quad Ocupado\\Algoritmos\\Prof. Danilo Rezende};
  \draw (0.3,5.3) rectangle (5.3,5.75);
  \node[anchor=west] at (0.4,5.52) {12:00 \quad Livre — toque para reservar};
  \draw (0.3,4.85) rectangle (5.3,5.3);
  \node[anchor=west] at (0.4,5.07) {13:00 \quad Livre — toque para reservar};
  \node at (2.8,4.62) {$\vdots$};
  \draw[fill=gray!30] (0.3,3.25) rectangle (5.3,4.4);
  \node[anchor=north west, text width=4.8cm] at (0.4,4.3) {\textbf{19:00 às 22:00} \quad Ocupado\\Banco de Dados\\Prof. Sérgio Santana};

  % legenda e botão
  \draw (0.3,2.75) rectangle (0.6,2.95); \node[anchor=west] at (0.65,2.85) {Livre};
  \draw[fill=gray!30] (1.7,2.75) rectangle (2.0,2.95); \node[anchor=west] at (2.05,2.85) {Ocupado};
  \draw[fill=black!80, rounded corners=3pt] (0.3,1.8) rectangle (5.3,2.4);
  \node[white] at (2.8,2.1) {\textbf{+ Nova reserva}};
  \node[text width=5cm, align=center, gray] at (2.8,1.0) {No computador, a mesma agenda aparece\\em grade com a semana inteira.};
\end{tikzpicture}

    \captionof{figure}{Tela 1 · Agenda do laboratório.}
  \end{minipage}\hfill
  \begin{minipage}[t]{0.48\linewidth}
    \centering
    % Tela 2 · Nova reserva (H02 · H01 · RF07)
\begin{tikzpicture}[font=\sffamily\scriptsize,
  campo/.style={draw, minimum height=0.5cm, anchor=north west, inner sep=3pt, text depth=0pt}]
  % aparelho e barra superior
  \draw[thick, rounded corners=8pt] (0,0) rectangle (5.6,10.4);
  \fill[gray!25] (0.2,9.6) rectangle (5.4,10.2);
  \node[anchor=west] at (0.3,9.9) {\textbf{Reserva de Laboratórios}};
  \node[anchor=east] at (5.3,9.9) {Sair};
  \node[anchor=west] at (0.3,9.25) {\small\textbf{Nova reserva}};

  % campos
  \node[anchor=west] at (0.3,8.85) {Professor};
  \draw[fill=gray!15] (0.3,8.1) rectangle (5.3,8.6);
  \node[anchor=west, gray] at (0.4,8.35) {Wilia F. da Silva (vem do login)};

  \node[anchor=west] at (0.3,7.85) {Laboratório};
  \draw (0.3,7.1) rectangle (5.3,7.6);
  \node[anchor=west] at (0.4,7.35) {Lab 101 (40 lugares)};
  \node[anchor=east] at (5.2,7.35) {$\blacktriangledown$};

  \node[anchor=west] at (0.3,6.85) {Data};
  \draw (0.3,6.1) rectangle (5.3,6.6);
  \node[anchor=west] at (0.4,6.35) {17/11/2026};

  \node[anchor=west] at (0.3,5.85) {Início};
  \draw (0.3,5.1) rectangle (2.6,5.6);
  \node[anchor=west] at (0.4,5.35) {13:00};
  \node[anchor=east] at (2.5,5.35) {$\blacktriangledown$};
  \node[anchor=west] at (3.0,5.85) {Duração};
  \draw (3.0,5.1) rectangle (5.3,5.6);
  \node[anchor=west] at (3.1,5.35) {2 horas};
  \node[anchor=east] at (5.2,5.35) {$\blacktriangledown$};
  \node[anchor=west] at (0.3,4.85) {Término: \textbf{15:00} (calculado)};

  \node[anchor=west] at (0.3,4.45) {Disciplina};
  \draw (0.3,3.7) rectangle (5.3,4.2);
  \node[anchor=west] at (0.4,3.95) {Engenharia de Software};

  \node[anchor=west] at (0.3,3.45) {Número de alunos};
  \draw (0.3,2.7) rectangle (2.6,3.2);
  \node[anchor=west] at (0.4,2.95) {30};
  \node[anchor=west, gray] at (2.7,2.95) {máx. 40 neste lab};

  % disponibilidade e botão
  \draw[fill=green!12, draw=green!50!black] (0.3,1.95) rectangle (5.3,2.45);
  \node[anchor=west] at (0.4,2.2) {$\checkmark$ Lab 101 livre das 13:00 às 15:00};
  \draw[fill=black!80, rounded corners=3pt] (0.3,1.1) rectangle (5.3,1.7);
  \node[white] at (2.8,1.4) {\textbf{Confirmar reserva}};
  \node[gray, text width=5cm, align=center] at (2.8,0.55) {Duração só de 1 a 4 horas.\\Horários fechados não aparecem.};
\end{tikzpicture}

    \captionof{figure}{Tela 2 · Nova reserva.}
  \end{minipage}
\end{center}

\begin{center}
  % Tela 3 · Minhas reservas (H03 · H04 · H06)
\begin{tikzpicture}[font=\sffamily\scriptsize]
  % aparelho e barra superior
  \draw[thick, rounded corners=8pt] (0,0) rectangle (5.6,10.4);
  \fill[gray!25] (0.2,9.6) rectangle (5.4,10.2);
  \node[anchor=west] at (0.3,9.9) {\textbf{Reserva de Laboratórios}};
  \node[anchor=east] at (5.3,9.9) {Sair};
  \node[anchor=west] at (0.3,9.25) {\small\textbf{Minhas reservas}};

  % aviso de confirmação
  \draw[fill=green!12, draw=green!50!black] (0.3,8.2) rectangle (5.3,8.95);
  \node[anchor=west, text width=4.8cm] at (0.4,8.57) {$\checkmark$ Reserva confirmada. Aviso enviado para o seu e-mail.};

  % reserva 1
  \draw[rounded corners=3pt] (0.3,6.4) rectangle (5.3,7.95);
  \node[anchor=north west] at (0.4,7.9) {\textbf{Lab 101} · Ter, 17/11 · 13:00 às 15:00};
  \node[anchor=north west] at (0.4,7.52) {Engenharia de Software · 30 alunos};
  \draw[fill=green!15, draw=green!50!black, rounded corners=4pt] (0.4,6.55) rectangle (2.0,6.85);
  \node at (1.2,6.7) {Reservado};
  \draw[rounded corners=2pt] (2.4,6.5) rectangle (3.7,6.9); \node at (3.05,6.7) {Editar};
  \draw[rounded corners=2pt] (3.9,6.5) rectangle (5.2,6.9); \node at (4.55,6.7) {Cancelar};

  % reserva 2
  \draw[rounded corners=3pt] (0.3,4.65) rectangle (5.3,6.2);
  \node[anchor=north west] at (0.4,6.15) {\textbf{Lab 202} · Qui, 19/11 · 19:00 às 21:00};
  \node[anchor=north west] at (0.4,5.77) {Banco de Dados · 25 alunos};
  \draw[fill=green!15, draw=green!50!black, rounded corners=4pt] (0.4,4.8) rectangle (2.0,5.1);
  \node at (1.2,4.95) {Reservado};
  \draw[rounded corners=2pt] (2.4,4.75) rectangle (3.7,5.15); \node at (3.05,4.95) {Editar};
  \draw[rounded corners=2pt] (3.9,4.75) rectangle (5.2,5.15); \node at (4.55,4.95) {Cancelar};

  % reserva cancelada
  \draw[rounded corners=3pt, draw=gray] (0.3,2.9) rectangle (5.3,4.45);
  \node[anchor=north west, gray] at (0.4,4.4) {\textbf{Lab 102} · Seg, 16/11 · 08:00 às 10:00};
  \node[anchor=north west, gray] at (0.4,4.02) {Algoritmos · 35 alunos};
  \draw[fill=red!12, draw=red!50!black, rounded corners=4pt] (0.4,3.05) rectangle (2.0,3.35);
  \node at (1.2,3.2) {Cancelada};

  \node[text width=5cm, align=center, gray] at (2.8,1.8) {Cancelar pede confirmação.\\Editar abre o formulário já preenchido\\e verifica de novo as regras.};
\end{tikzpicture}

  \captionof{figure}{Tela 3 · Minhas reservas.}
\end{center}

% ---------------------------------------------------------------------
\section{Artigo: introdução}

\textbf{Tema 6 · Teste de usabilidade em aplicativos}

Um sistema pode cumprir todas as funções previstas e, ainda assim,
fracassar quando as pessoas que o usam não entendem o que fazer,
cometem erros com frequência ou desistem da tarefa. Por isso, além de
verificar se o software funciona, é preciso verificar se ele é fácil
de usar. Segundo Barbosa e Silva (2010), a usabilidade está relacionada
à facilidade de aprendizado e de uso de um sistema e à satisfação de
quem o utiliza, e deve ser avaliada com métodos próprios ao longo do
desenvolvimento.

Esse cuidado vale também para sistemas pequenos, criados para resolver
problemas do dia a dia. É o caso do sistema de reserva de laboratórios
de uma faculdade: a coordenação recebia os pedidos por e-mail e era
comum dois professores chegarem ao mesmo laboratório no mesmo horário.
A versão 1.0 do sistema foi criada para mostrar as reservas e impedir
esses choques. No entanto, se os professores tiverem dificuldade para
reservar, a tendência é que voltem a usar o e-mail, e o problema
original continua.

Avaliar a usabilidade não exige, necessariamente, muitos participantes.
O modelo proposto por Nielsen e Landauer (1993) mostra que a quantidade
de problemas encontrados cresce rapidamente com os primeiros usuários
testados e que um grupo pequeno, por volta de cinco pessoas, já revela
a maior parte deles. A observação dos participantes pode ser combinada
com um questionário de satisfação, como a escala SUS
(BROOKE, 1996), e com a classificação dos problemas pelas heurísticas
de usabilidade (NIELSEN; MOLICH, 1990).

Diante disso, este artigo busca responder à pergunta: \emph{como testar
se um aplicativo é fácil de usar?} Para isso, apresenta um estudo de
caso em que a versão 1.0 do sistema de reserva de laboratórios foi
testada por cinco colegas de curso. Os objetivos específicos são:
(i) planejar tarefas que representem o uso real do sistema;
(ii) medir a taxa de sucesso e o tempo de cada tarefa;
(iii) aplicar o questionário SUS; e
(iv) classificar os problemas encontrados e propor melhorias para a
versão 2.0.

O restante do texto está organizado da seguinte forma: a Seção 2
apresenta a fundamentação teórica sobre usabilidade e testes com
usuários; a Seção 3 descreve o desenvolvimento do estudo e analisa os
resultados; e a Seção 4 traz as considerações finais.

\subsection*{Referências citadas na introdução}

{\small
BARBOSA, S. D. J.; SILVA, B. S. \emph{Interação humano-computador}. Rio de Janeiro: Elsevier, 2010.

BROOKE, J. SUS: a ``quick and dirty'' usability scale. In: JORDAN, P. W. et al. (Ed.). \emph{Usability Evaluation in Industry}. London: Taylor \& Francis, 1996. p. 189--194.

NIELSEN, J.; LANDAUER, T. K. A mathematical model of the finding of usability problems. In: \emph{Proceedings of the INTERCHI'93 Conference on Human Factors in Computing Systems}. Amsterdam: ACM, 1993. p. 206--213. DOI: 10.1145/169059.169166.

NIELSEN, J.; MOLICH, R. Heuristic evaluation of user interfaces. In: \emph{Proceedings of the SIGCHI Conference on Human Factors in Computing Systems (CHI'90)}. Seattle: ACM, 1990. p. 249--256. DOI: 10.1145/97243.97281.
}

\end{document}
